\documentclass[UTF8,a4paper,11pt,openany]{ctexbook}
\usepackage{../common/statlearnbook}
\SLSetBookMetadata{第 1 册}{统计学习方法讲义 v2.6.0}{数学基础与统计学习基本理论}
\begin{document}
\setcounter{page}{182}
\setcounter{chapter}{10}
\setcounter{figure}{0}
\pagestyle{headings}
\chaptermark{模型评估、选择与泛化}
\begin{optimizationbox}
\textbf{优化解释。}$\ell_2$惩罚提高光滑目标的曲率并缓解病态；$\ell_1$目标在零点不可微，需用次梯度、近端或坐标方法。正则化强度既改变统计偏差与方差，也改变数值条件，不能只按优化收敛速度选择。
\end{optimizationbox}
\begingroup
\tikzset{every picture/.append style={scale=1.12}}
% v2.3.1 figure UID: FIG-P157-01
% Proposed post-v2.3.1 polish for FIG-P157-01: protect key-label size and stroke hierarchy.
\tikzset{
  slfig-FIG-P157-01/.style={every path/.append style={line cap=round,line join=round}},
  slfig-FIG-P157-01-direct/.style={slfig direct label,
    font=\fontsize{9.2pt}{10.9pt}\selectfont,
    fill=white,fill opacity=.90,text opacity=1,inner sep=.9pt},
  slfig-FIG-P157-01-key/.style={font=\fontsize{9.2pt}{10.9pt}\selectfont,
    fill=white,fill opacity=.90,text opacity=1,inner sep=.8pt},
  slfig-FIG-P157-01-region/.style={font=\fontsize{8.8pt}{10.4pt}\selectfont}
}
\pgfplotsset{slfig-FIG-P157-01-axis/.style={
  tick label style={font=\fontsize{8.6pt}{10.2pt}\selectfont},
  label style={font=\fontsize{9.4pt}{11.1pt}\selectfont},
  xlabel style={at={(axis description cs:.5,-.18)},anchor=north},
  clip=false}}
\begin{figure}[H]
\centering
\begin{tikzpicture}[slfig-FIG-P157-01]
\begin{axis}[slfig-FIG-P157-01-axis,
  slfig axis,
  width=.85\linewidth,height=.47\linewidth,
  xmin=0,xmax=10,ymin=0,ymax=4.35,
  axis lines=left,
  xlabel={模型复杂度},ylabel={预测误差},
  xtick=\empty,ytick=\empty,
  axis on top,
  clip=false,
]
  \path[fill=SLBlue!3] (axis cs:0,0) rectangle (axis cs:3.65,4.35);
  \path[fill=SLGold!5] (axis cs:3.65,0) rectangle (axis cs:6.70,4.35);
  \path[fill=SLRed!3] (axis cs:6.70,0) rectangle (axis cs:10,4.35);
  \addplot[draw=SLBlue,line width=1.00pt,domain=0:10,samples=241]
    {0.36+3.35*exp(-0.34*x)};
  \addplot[draw=SLTeal,line width=1.05pt,dash pattern=on 4pt off 2.4pt,
    domain=0:10,samples=241]
    {1.08+0.105*(x-5.25)^2};
  \addplot[slfig reference,line width=.70pt,dash pattern=on 3pt off 2pt]
    coordinates {(5.25,0) (5.25,1.08)};
  \addplot[only marks,mark=*,mark size=2.4pt,draw=SLGold,fill=SLGold]
    coordinates {(5.25,1.08)};
  \draw[draw=SLBlue,line width=.48pt]
    (axis cs:7.15,.655)--(axis cs:7.35,.90);
  \node[slfig-FIG-P157-01-direct,fill=none,text=SLBlue,anchor=west]
    at (axis cs:7.43,1.02)
    {训练误差：单调下降};
  \node[slfig-FIG-P157-01-direct,text=SLTeal,anchor=west] at (axis cs:3.9,2.2)
    {验证误差：先降后升};
  \node[slfig-FIG-P157-01-key,text=SLGold,anchor=south] at (axis cs:5.25,1.18)
    {最低验证误差};
  \node[slfig-FIG-P157-01-key,anchor=north] at (axis cs:5.25,-.07)
    {选择复杂度};
  \node[slfig-FIG-P157-01-region,text=SLGray,anchor=north]
    at (axis cs:1.82,-.44) {欠拟合};
  \node[slfig-FIG-P157-01-region,text=SLGold,anchor=north]
    at (axis cs:5.18,-.44) {合适};
  \node[slfig-FIG-P157-01-region,text=SLRed,anchor=north]
    at (axis cs:8.35,-.44) {过拟合};
\end{axis}
\end{tikzpicture}
\caption{模型复杂度增加时训练误差通常下降，而验证误差可能先降后升。}\label{fig:V1-C10-complexity}
\end{figure}

\endgroup
在\cref{fig:V1-C10-complexity}上选择容量时只使用验证曲线和预先约定的规则；训练误差持续下降并不构成增大容量的证据。候选锁定后还需在未参与选择的数据上报告一次最终误差。
\end{document}

